% !TeX program = xelatex
% 第二版语言修订稿；2026-10-25；单文件，参考文献及结果内嵌。
% 构造聚类研究：scripts/semantic_clustering_study.py。
% 真实分类案例重分析：scripts/title_case_analysis.py。
\documentclass[UTF8,a4paper,zihao=-4,fontset=fandol]{ctexart}
\usepackage{fontspec}
\setmainfont{Latin Modern Roman}
\usepackage[margin=25mm,headheight=15pt]{geometry}
\usepackage{amsmath,amssymb,amsthm,booktabs,array,tabularx}
\usepackage[authoryear,round]{natbib}
\usepackage{xcolor,fancyhdr,enumitem}
\usepackage[section]{placeins}
\usepackage{xurl}
\usepackage[unicode,colorlinks=true,linkcolor=black,citecolor=black,urlcolor=blue!45!black]{hyperref}
\hypersetup{pdftitle={语义决策模型在统计分析中的应用：方法论框架与研究议程},pdfauthor={游尹哲郡},pdfkeywords={语义决策模型,文本分类,约束聚类,序数测量,记录链接}}
\urlstyle{same}
\linespread{1.10}
\setlength{\parindent}{2em}
\setlength{\parskip}{0.12em}
\setlength{\emergencystretch}{2em}
\setlength{\tabcolsep}{4pt}
\renewcommand{\arraystretch}{1.16}
\newcolumntype{Y}{>{\raggedright\arraybackslash}X}
\newcolumntype{L}[1]{>{\raggedright\arraybackslash}p{#1}}
\pagestyle{fancy}
\fancyhf{}
\fancyhead[L]{\small 游尹哲郡}
\fancyhead[R]{\small 语义决策模型与统计分析}
\fancyfoot[C]{\small\thepage}
\renewcommand{\headrulewidth}{0.3pt}
\setlist[enumerate]{leftmargin=2em,itemsep=0.25em,topsep=0.4em}
\newcommand{\E}{\mathbb E}
\newcommand{\Prob}{\mathbb P}
\newcommand{\Var}{\operatorname{Var}}
\newcommand{\ind}{\mathbf 1}
\DeclareMathOperator*{\argmin}{arg\,min}
\newtheorem{proposition}{命题}
\title{\bfseries 语义决策模型在统计分析中的应用：\\方法论框架与研究议程}
\author{游尹哲郡（Youyin Zhejun）\\[0.3em]
\small 中央财经大学\\
\small ORCID：\href{https://orcid.org/0009-0009-5316-8193}{0009-0009-5316-8193}}
\date{2026年10月25日}

\begin{document}
\maketitle
\thispagestyle{plain}
\begin{abstract}
语义决策模型依据自然语言标准生成类别、关系和等级判断，可供统计分析构造监督信息、分组结构与测量变量。本文结合定向文献研究、真实分类案例重分析与构造数值研究，建立涵盖研究标准、判断获取、概率读取和整体结构求解的方法论框架。论文标题提取错误识别案例使用同一批1,000条测试记录；在固定基础权重、NF4 量化、模板与评分程序的配对比较中，Tev1 任务微调后的宏平均 F1 从71.13\%提高至89.69\%，漏检由242条降至34条，误报由7条增至68条。非空记录的 Brier 分数由0.1685降至0.0993，分组后的预测概率与参考频率仍有差异。聚类研究对九条构造记录的全部划分精确求解，发现概率加权的表现受判断可靠度、查询覆盖及并列最优解选择规则共同影响；完整图上的方法排序在替代规则下反转。序数测量与记录链接分别讨论等级分布的测量含义和实体关系对统计单位的影响。分类误差、概率质量、结构一致性与下游统计结果据此分别评价，后续比较按标签预算、部署条件和总成本组织。
\end{abstract}
\noindent\textbf{关键词：}语义决策模型；监督信息获取；约束聚类；概率加权；序数测量；记录链接。

\section{引言}\label{sec:intro}

统计分析需要先从观测记录中识别与研究问题相关的信息。客户投诉可以按产品类别编码，也可以按故障机制分组；访谈材料可以记录某项态度是否出现，也可以依据评分规程构造程度变量。类别、关系和等级的获取方式决定了后续统计方法接收的信息。这些信息可由人工标注、已有训练标签或表示空间中的距离获得。研究标准发生变化时，重新编码和重新组织资料通常需要额外投入。

近年来，自然语言推断模型、通用语言模型及专门决策模型均可依据文字说明产生判断。零样本分类利用标签描述定义目标任务；语言模型辅助标注将编码规程纳入查询；Jev 等专门接口进一步规定答案类型及返回字段\citep{yin2019,tan2024,typesafe-intro}。由此形成一个统计方法论问题：研究者能够批量获取类别判断、成对关系与等级概率后，监督信息、分组结构和测量变量应如何构造？这些信息在何种任务和预算条件下改善统计分析？

本文以\emph{语义决策模型}统称依据明确语义标准提供判断的模型，以\emph{决策接口}指称研究者实际调用的操作。Jev 与 Tev1 是本文讨论的两种专门实现，真实分类案例使用 Tev1。分析围绕语义标准、输入记录、答案空间与输出字段的组合展开。通用语言模型、零样本分类模型和专门模型虽有实现差异，也可以执行相同判断任务，因此可按统计用途比较其提供的信息。

已有方法为此提供了明确基础。弱监督学习利用带噪标签来源训练预测器\citep{ratner2017}；约束聚类接收同组与异组信息\citep{wagstaff2001}；ClusterLLM 和少样本聚类研究利用语言模型判断指导表示、约束或聚类粒度\citep{clusterllm2023,viswanathan2024}。Cequel 进一步研究有限查询预算下的信息选择和加权约束，近期研究也采用置信度阈值与惩罚机制处理自动生成的约束\citep{cequel2025,jia2026}。因此，“新范式”应接受逐项考察：新增的能力发生在哪个信息获取环节，与哪些既有方法衔接，其收益来自何种信息及处理方式。

\citet{tan2024}围绕标注生成、评估与使用组织综述，ClusterLLM、少样本聚类及 Cequel 针对聚类提出具体方法。本文从信息获取条件及统计方法的假设出发，统一讨论标签、关系和等级在不同任务中的用途。真实分类案例分析任务监督、部署配置与错误代价；构造聚类比较保持关系判断相同，考察概率读取与整体划分求解的作用。本文沿用既有统计与学习方法中的目标函数、误差分析和估计工具，贡献在于明确应用选择、比较条件及各类证据所回答的问题。

本文开展定向叙述性文献研究。既有方法与 Jev 材料核查截至2026年10月5日，Tev1 与分类案例材料补充核查至10月6日。方法文献主要取自 ACL Anthology、会议论文集、出版方记录及作者原文；模型材料包括官方文档、固定版本模型卡与可核查的研究论文。官方资料用于确认操作和字段，经验研究用于评价规定任务下的表现。标题案例复用2026年10月3日至6日完成的课程实验及逐条预测，本文新增指标复算、配对重抽样与错误代价分析。构造聚类研究分析明确设定下的统计后果；既有 SMS 抽样结果作为下游补充案例。

全文依次说明模型形态与方法前史，建立信息用途框架，展开分类与聚类两条主线，再讨论序数测量与记录链接。第\ref{sec:title-case}节报告真实分类案例，第\ref{sec:evaluation}节报告构造聚类结果及比较评价方案，附录\ref{sec:sms}保留 SMS 下游应用，第\ref{sec:conclusion}节总结适用条件。

\section{模型形态、接口属性与方法前史}\label{sec:history}

\subsection{结构化判断提供的操作}

Jev 文档规定三类判断操作。Choice 接收候选选项及说明，返回选项、各选项概率和置信度；Score 接收有序等级及其描述，返回等级概率、索引加权得分和置信度；Noul 返回对给定问题回答“是”的赋值概率\citep{typesafe-choice,typesafe-score,typesafe-noul}。这些字段分别适合表达名义类别、序数等级和二元关系。文档还允许同一输入接收多项问题\citep{typesafe-intro}，便于在同一批记录上获取多种研究属性。

Tev1 是基于 Qwen3.5-4B 监督微调的实验性决策模型。其输入包括状态、问题和带字母标识的选项，目标输出为选项字母，应用程序据此确定语义标签\citep{tev1-card}。真实案例的 NF4 实现在答案位置读取 A、B 两个词元的 logit，归一化后得到二类分数。结构化判断既可由服务接口返回，也可由本机推理程序组织，概率字段的定义应随具体读取程序说明。

Choice 的置信度由返回分布的集中程度计算，Score 的数值得分由等级索引及概率计算\citep{typesafe-choice,typesafe-score}。概率校准、等级效度与跨任务可比性分别依据参考标签、评分规程和应用总体评价。同一请求中的多项判断共享文本和模型，问题标准也可能接近，其统计相关性应在数据分析中估计。

决策接口允许研究者在判断操作中直接指定自然语言标准，并通过同一类请求获取类别、关系和等级。返回的完整概率分布可进入损失函数或训练目标，同一批记录也可依据多个标准重新评价。模型训练、目标任务核验和统计优化的投入应分别记录；速度、价格与准确性的比较还应说明模型版本、任务范围、请求规模及总成本。

\subsection{既有方法中的对应位置}

监督分类依赖特征与标签的对应关系，半监督学习进一步利用未标记记录的结构。自然语言推断式零样本分类使用记录与类别描述之间的蕴含关系确定标签\citep{yin2019}；弱监督框架将多种带噪判断建模并用于训练\citep{ratner2017}。语义决策输出可以承担直接编码或弱标签来源，其质量与获取成本由目标任务验证。

聚类方法利用表示、距离、密度或概率模型定义分组结构。约束聚类增加外部关系信息，语言模型辅助聚类则研究这些信息如何获取。ClusterLLM 使用三元组判断指导表示学习、成对判断选择粒度；\citet{viswanathan2024}比较特征改进、约束和结果修正；Cequel 研究信息对和三元组的选择\citep{clusterllm2023,cequel2025}。专门决策接口可以提供关系来源，但所请求的关系、算法接收的约束与评价目标须对应。

\begin{table}[htbp]
\centering\small
\caption{语义决策操作与既有方法的联系}\label{tab:history}
\begin{tabularx}{\textwidth}{@{}L{2.1cm}YY@{}}
\toprule
操作环节 & 既有基础 & 专门接口的可评价变化 \\
\midrule
表达类别标准 & 标签描述、自然语言推断、编码规程 & 类别说明的组织方式；标准修订后的编码工作量 \\
获取关系信息 & 人工约束、成对相似性、三元组偏好 & 固定标准下的关系误差、查询数与费用 \\
读取概率输出 & 概率分类、软标签、加权损失 & 完整分布的质量；读取与校准的成本 \\
复用观测记录 & 多属性标注、多任务学习 & 多标准评价的一致性与核验负担 \\
\bottomrule
\end{tabularx}
\end{table}

表\ref{tab:history}所列变化均有对应的方法基础。不同模型接受相同标准并返回同类判断时，可在固定统计方法下比较判断质量与获取预算；同一模型的输出则可用于比较不同统计使用方式。

\section{语义判断进入统计分析的统一框架}\label{sec:framework}

\subsection{从研究标准到分析对象}

设 $x_i$ 为第 $i$ 条记录，$q$ 为研究者制定的判断标准，$s(x_i)$ 为提供给模型的信息，$\mathcal A$ 为答案空间，$m$ 为固定模型及版本。一个判断程序记为
\begin{equation}
 O_i=M_m\{s(x_i);q,\mathcal A\},\qquad
 \widehat\theta=\mathcal T\bigl(\{O_i\},\{Y_i^{\mathrm{ref}}:i\in V\};\eta\bigr).
 \label{eq:framework}
\end{equation}
其中 $O_i$ 为结构化输出，$V$ 为核验记录集合，$Y_i^{\mathrm{ref}}$ 为依规定程序获得的参考结果，$\mathcal T$ 为统计方法，$\eta$ 为其参数及选择规则。成对关系的输入相应为 $(x_i,x_j)$，三元组判断的输入为 $(x_i,x_j,x_k)$。式\eqref{eq:framework}明确了模型提供信息、统计方法使用信息的接口。

研究概念规定所关心的属性，判断规程说明识别该属性所需的证据，参考结果记录规程的实施。例如，评价“严重投诉”时，研究者须说明严重程度的含义，制定等级和证据要求，再由核验人员依规程评分。模型与参考结果的一致性评价编码是否依规程执行，等级与外部事实的联系用于评价测量效度。类别定义、上下文或等级描述发生修订后，评价材料应对应修订后的判断程序。文本分析中的实质性解释与任务验证已有系统讨论\citep{grimmer2013}，本文据此考察输入标准及其统计用途。

\subsection{标签、关系与等级的统计用途}

\begin{table}[htbp]
\centering\small
\caption{语义信息进入统计方法的主要路径}\label{tab:framework}
\begin{tabularx}{\textwidth}{@{}L{1.7cm}L{2.7cm}YY@{}}
\toprule
统计任务 & 获取的信息 & 进入方法的方式 & 主要评价对象 \\
\midrule
分类与标注 & 类别判断、类别概率、多属性编码 & 直接编码；弱监督融合；学生模型训练；人工复核 & 记录误差、稀有类别、核验量、标准修订成本 \\
聚类 & 同组关系、成对比较、三元组偏好 & 约束聚类；关系图；距离或表示学习 & 目标划分、图覆盖、冲突、概率权重、查询预算 \\
序数测量 & 按规程获得的等级分布 & 序数变量；阈值指标；多维测量 & 评分一致性、等级使用、跨组可比性、编码敏感性 \\
记录链接 & 同一实体判断、比较证据 & 候选匹配；概率链接；整体匹配 & 候选覆盖、错误合并、实体计数、链接后分析 \\
\bottomrule
\end{tabularx}
\end{table}

表\ref{tab:framework}中的输出类型确定了方法可以利用的信息。名义标签用于类别学习，成对关系用于组织整体结构，有序等级用于测量。记录链接同样接收关系信息，其关系目标为实体同一性。四类任务均可依次考察传统方法所需的信息、模型获取方式、统计使用规则及收益验证。

信息质量的评价涉及局部判断、整体结构和下游分析。记录或记录对与参考结果的比较反映局部误差；划分、等级次序和实体映射反映结构；回归、总体汇总或决策结果反映下游分析后果。即使局部误差率相同，错误所在位置也可能影响分析结果。例如，跨组桥接可能影响聚类划分，极端等级遗漏会影响序数测量，错误合并则会改变链接后的实体构成。

\subsection{预算与优势的评价条件}

语义判断的收益应结合任务误差与总投入评价。记标准制定成本为 $C_q$，模型判断数为 $B$，每次平均成本为 $c_M$，参考标签数为 $v$，平均标签成本为 $c_H$，训练及后续处理成本为 $C_T$，则
\begin{equation}
 C=C_q+B c_M+v c_H+C_T.
 \label{eq:cost}
\end{equation}
一次接口请求可以包含多项判断，$B$ 按语义问题计数；请求数、输入输出规模与实际费用分别报告。训练、验证和最终核验的标签分项记录，各阶段单位成本不同时分别加总；已有标签的获取成本与本次新增投入也分别说明。对于反复变更分类标准的项目，$C_q$ 及新标准的核验投入应累计。

目标任务标签较少、判断标准清楚且概率输出可重复使用时，语义判断值得纳入方法比较。其优势可按同一费用下的目标误差，或达到同一误差要求所需的总投入评价。比较应采用相同的类别或关系目标，最终评价材料应独立于标准调整过程。

\section{分类：监督信息的获取与组织}\label{sec:classification}

\subsection{直接编码与目标任务标签}

分类任务具有事先规定的类别清单和归类标准。传统训练流程通过标记样本估计特征与类别的对应关系；语义判断允许类别说明直接参与预测。零样本设定指目标任务训练标签数量为零，预训练资料仍为模型提供语言及领域信息\citep{yin2019}。这一信息来源差异应在解释标签预算时说明。

设模型返回 $K$ 个互斥类别的概率 $\widehat p_{ik}$，直接编码可以选择最大概率类别，也可以依据误分类成本选择行动：
\begin{equation}
 \widehat a_i=\argmin_{a\in\mathcal D}\sum_{k=1}^K L(a,k)\widehat p_{ik}.
 \label{eq:classification}
\end{equation}
行动空间 $\mathcal D$ 可包含人工复核，损失 $L$ 相应包含漏检、误报和复核成本。该规则的评价需要核验概率质量及成本设定。类别中设置“证据不足”等选项时，其含义须与记录内容和研究目标一致。

同一记录涉及多项属性时，可以分别判断产品、故障、诉求和程度。每项属性遵循自身标准，多项判断的联合分布受共同文本影响。单个互斥分类问题中的概率分配描述模型对候选类别的判断；多个属性问题的联合结果记录属性共现。二者对应不同的数据组织方式。

\subsection{任务适配、弱监督与学生模型训练}

决策模型可在给定语义规程下继续使用目标任务标签训练，参数估计同时利用规程描述与参考编码。低秩适配在冻结主干的条件下学习附加参数，QLoRA 采用量化主干完成这种训练\citep{qlora2023}。评价时应分别记录训练标签、验证标签及训练投入，并设置部署条件相同的零样本对照。第\ref{sec:title-case}节采用这一设计。

语义决策输出可以与规则、词典和已有分类器共同构成弱标签来源。弱监督方法估计来源质量及相关性，再形成训练标签\citep{ratner2017}。多个提示或多个问题来自同一模型时，共同训练信息和相近标准可能产生共有错误。来源数量与新增信息量的关系需要在参考样本上评价。

完整类别分布也可以监督较小模型。令学生模型输出 $f_{\theta k}(x_i)$，软标签训练采用
\begin{equation}
 \widehat\theta=\argmin_\theta
 -\sum_{i=1}^n\sum_{k=1}^K\widehat p_{ik}\log f_{\theta k}(x_i).
 \label{eq:student}
\end{equation}
该目标保留教师模型在类别间的概率分配，适合研究一次语义编码能否支持后续批量预测。实验应在独立参考标签上比较硬标签训练、软标签训练和相同人工标签预算下的模型，并计入教师查询与学生训练成本。学生与教师的一致性可用于评价知识传递，学生与参考标签的误差用于评价目标任务表现。

\subsection{人工复核与分类评价}

人工复核可以优先检查低置信记录、稀有类别和错误代价较高的记录。抽样方案随复核目标而定：修订困难记录、估计分类误差与估计总体类别比例，对样本覆盖的要求各有差异。用于总体评价的核验样本应记录选择概率，并覆盖研究总体中的相关记录类型。

分类比较宜报告宏平均 F1、类别召回率、混淆矩阵和概率评分，并考察文本来源或群体之间的差异。Brier 分数评价概率与二元参考结果的平方误差；校准分析评价报告概率与实际频率的关系\citep{guo2017}。研究标准调整后，还应计算修订规程、重查记录和重新核验的工作量。语言模型标注综述已经覆盖生成、评估和使用\citep{tan2024}，专门接口的应用研究可进一步比较上述工作量与规定字段的使用效果。

\citet{ibrahim2026}在计算社会科学任务中比较决策模型与语言模型，报告了所研究决策模型较低的查询费用、随任务变化的分类质量，并评价置信度路由。其结果为按任务选择判断来源和复核方式提供了经验依据。第\ref{sec:agenda}节据此分别控制判断来源与使用方式设计比较。

\subsection{真实案例：论文标题提取错误识别}\label{sec:title-case}

\subsubsection{数据、语义标准与比较设计}

本案例复用信息检索课程项目“论文标题提取错误识别”的冻结材料\citep{title-case}。目标为判断机器提取文本是否属于错误提取，参考结果为课程文件提供的标签。原训练文本共236,158行，依次清理空行、类内重复、测试标题重叠及训练标签冲突，形成212,385条记录的训练池。标题身份按 Unicode NFKC、大小写折叠、空白合并及首尾清理定义；种子42的分层划分产生169,908条训练记录与42,477条验证记录。测试工作簿包含568条正确提取、432条错误提取，共1,000条记录。本次重分析独立读取处理档案，确认三个分区的归一化标题交集均为空。

模型输入为机器提取标题；人工原题名、文献 ID 与参考标签供结果核对。规程将完整、连贯的学术标题归为 valid，将残缺片段、作者或机构名单、出版元数据、页眉、乱码及混入正文的文本归为 invalid。两个标签分别对应 A 与 B，输入结构为状态、问题及选项说明。测试集的24条空标题均有错误提取标签，各方法统一依据空标题规则判为 invalid；其余976条记录由模型判断。

比较包括词频多项式朴素贝叶斯（NB）、平均 Word2Vec 向量加 SVM、冻结 BERT 表示加 SVM，以及 Tev1 的两个零样本部署与任务微调。前三种方法和微调方法共用训练与验证分区，参数或检查点按验证集宏平均 F1 选择。Tev1 的 NF4 零样本与微调结果在同一模型实例中关闭或启用任务适配器，固定基础权重版本、量化、聊天模板、选项顺序及逐条评分程序。任务微调采用秩8的 QLoRA，训练一轮；第5,310与10,620步的验证结果分别为85.45\%和85.91\%，后者在测试评价前入选。完整配置见附录\ref{app:materials}。

Ollama 零样本采用 GGUF 权重与 Q4\_K\_M 量化，通过决策接口返回选项分布；NF4 方案采用 Transformers 推理，在首个答案位置计算
\begin{equation}
 \widehat p_i=\frac{\exp(\ell_{iB})}{\exp(\ell_{iA})+\exp(\ell_{iB})},\qquad
 \widehat Y_i=\ind\{\widehat p_i>0.5\}.
 \label{eq:title-probability}
\end{equation}
其中 $Y_i=1$ 表示错误提取。两类零样本部署共用测试记录、任务说明及空标题规则，其量化、模板构造和评分实现同时变化。Ollama 结果用于说明完整部署的表现，NF4 配对结果用于评价任务适配。

\subsubsection{分类结果、概率质量与错误代价}

\begin{table}[htbp]
\centering\small
\caption{同一测试集上的标题提取错误识别（$n=1{,}000$）}\label{tab:title-results}
\begin{tabularx}{\textwidth}{@{}Yrrrrrr@{}}
\toprule
方法 & 训练标签 & 准确率 & 宏 F1 & 错误类召回 & 误报 & 漏检 \\
\midrule
词频＋NB & 169,908 & 74.70 & 74.19 & 70.14 & 124 & 129 \\
Word2Vec＋SVM & 169,908 & 71.10 & 70.96 & 74.31 & 178 & 111 \\
冻结 BERT＋SVM & 169,908 & 81.40 & 81.35 & 88.19 & 135 & 51 \\
Tev1/Q4\_K\_M 零样本 & 0 & 67.50 & 58.83 & 25.00 & 1 & 324 \\
Tev1/NF4 零样本 & 0 & 75.10 & 71.13 & 43.98 & 7 & 242 \\
Tev1/NF4＋QLoRA & 169,908 & 89.80 & 89.69 & 92.13 & 68 & 34 \\
\bottomrule
\end{tabularx}
\vspace{0.35em}
\parbox{\textwidth}{\footnotesize 准确率、宏平均 F1 和召回率以百分数表示。误报指正确提取被判为错误，漏检指错误提取被判为正确。训练标签指本任务用于拟合参数的记录数，监督方案另使用42,477条验证标签；预训练投入另行记录。各行包含同一24条空标题规则结果。表中为既有预测的独立复算。}
\end{table}

按表\ref{tab:title-results}的评价口径，NF4 微调方案的宏平均 F1 较配对零样本方案高18.57个百分点，错误类精确率从96.45\%降至85.41\%。逐条比较发现，微调纠正210条原误判，同时新增63条误判，准确率净增加14.70个百分点。本文以963个归一化非空标题组为单位，按参考类别分层进行10,000次配对自助重抽样；每次保留组内记录，并加回24条固定规则记录。宏平均 F1 差异的95\%百分位区间为15.46至21.78个百分点，描述固定训练结果与测试来源构成下的重抽样波动。相较冻结 BERT＋SVM，微调方案的宏平均 F1 高8.34个百分点，相应区间为6.08至10.61个百分点；误报与漏检分别少67条和17条。

概率评价采用976条非空记录，空标题的确定性规则结果另行处理。表\ref{tab:title-probability}显示，NF4 配对比较中的 Brier 分数、对数损失及分组校准误差均下降。微调方案在377条记录上给出的标题提取错误概率至少为0.9；这些记录的平均预测概率为0.9934，参考标签中的错误提取频率为0.9204。该分组中的概率赋值高于参考频率。用于后续复核或决策时，校准方法和阈值应结合应用代价选择。

\begin{table}[htbp]
\centering\small
\caption{Tev1 错误类概率的评价（非空记录，$n=976$）}\label{tab:title-probability}
\begin{tabularx}{\textwidth}{@{}Yrrr@{}}
\toprule
方案 & Brier 分数 & 对数损失 & 分组校准误差 \\
\midrule
Ollama/Q4\_K\_M 零样本 & 0.2181 & 0.6307 & 0.2349 \\
Transformers/NF4 零样本 & 0.1685 & 0.5111 & 0.1704 \\
Transformers/NF4＋QLoRA & 0.0993 & 0.3988 & 0.1301 \\
\bottomrule
\end{tabularx}
\vspace{0.35em}
\parbox{\textwidth}{\footnotesize 分组校准误差按错误类概率的10个等宽区间计算，各区间的平均概率与参考频率之差取绝对值，再按记录比例加权。它评价此项指定分组；Brier 与对数损失评价逐条概率。所有值由保存的预测复算，数值越低表示对应评分较好。}
\end{table}

宏平均 F1 汇总类别识别表现，应用损失则取决于两类错误的代价。设漏检和误报的单位代价为 $c_{\mathrm{FN}}>0$ 与 $c_{\mathrm{FP}}>0$，该测试集上 NF4 微调相对零样本的总误判损失变化为
\begin{equation}
 \Delta L=-208c_{\mathrm{FN}}+61c_{\mathrm{FP}}.
 \label{eq:title-cost}
\end{equation}
在固定预测下，当 $c_{\mathrm{FN}}/c_{\mathrm{FP}}>61/208\approx0.2933$ 时，微调方案的测试误判损失较低。实际比较还应计入标签、训练及核验成本；方案选择随错误代价和这些投入共同变化。

\subsubsection{语义标准、任务监督与适用范围}

零样本阶段依据自然语言标准定义归类任务，微调阶段利用任务标签调整判断边界，保存的选项概率可供概率评分和复核设计使用。这三个环节分别对应标准表达、监督信息获取与概率使用。NF4 配对结果显示本次任务适配提高了宏平均 F1，表\ref{tab:title-results}则报告各完整方案在该测试集上的表现。较少标签下的效率可通过标签预算曲线评价，专门决策训练相对通用基础模型的贡献可通过相同基础权重的对照研究评价。

课程参考标签评价实际提取结果，模型查询侧重所给文本的标题完整性。测试记录“Notation as a Tool of”是残缺片段，NF4 零样本判为 valid，微调判为 invalid；“Editorial Statement”具有完整题名的形式，但其原题名为“Special Issue: The Future of Biofeedback Instrumentation”，参考标签及微调判断均为 invalid。这两例为解释错误而选取，分别涉及文本形式和文本与原文的对应关系。后续可在同一规程下比较仅提供标题文本与增加原文信息时的判断，并记录参考编码分歧及裁决依据。

\section{聚类：研究标准指导的关系获取}\label{sec:clustering}

\subsection{关系目标与分组结构}

聚类关注记录之间的分组结构。给定表示 $z_i$，$K$-means 通过最小化组内平方距离形成划分；约束聚类进一步接收同组和异组信息\citep{wagstaff2001}。语义判断可以依据明确标准询问记录对的关系，例如“这两条投诉是否涉及相同的主要故障机制”。关系标准规定分组依据，类别名称和完整类别清单可以在划分完成后形成。若事先规定全部类别并逐条归类，则按照第\ref{sec:classification}节的分类设计评价。

产品名称和故障描述可能同时影响文本表示。研究目标若为故障机制，而所用距离对产品差异赋予较大权重，直接询问故障关系可以补充与目标相关的信息。在固定查询预算与关系误差下，这些信息能否改善机制划分，是该例需要评价的问题。表示学习和聚类约束均可利用此类反馈，已有实现包括 ClusterLLM、少样本聚类和 Cequel\citep{clusterllm2023,viswanathan2024,cequel2025}。

关系标准决定适用的分组结构。“共享任一故障因素”允许一条记录与多组存在关系；“主要故障机制相同”在主要机制唯一的规程下定义互斥分组。多机制记录适合采用关系图或重叠聚类，互斥划分则需要明确主要机制的选择规则。

\subsection{查询图、对称性与冲突}

设 $E$ 为获得判断的记录对集合，$R_{ij}^{(q)}$ 为按标准 $q$ 定义的同组关系，模型返回 $\widehat p_{ij}$。同组关系具有对称性。对于分别查询的两个输入顺序，可采用 $\bar p_{ij}=(\widehat p_{ij}+\widehat p_{ji})/2$，并记录顺序差异 $|\widehat p_{ij}-\widehat p_{ji}|$。平均值的方差为
\begin{equation}
 \Var(\bar p_{ij})=
 \frac{\Var(\widehat p_{ij})+\Var(\widehat p_{ji})
 +2\operatorname{Cov}(\widehat p_{ij},\widehat p_{ji})}{4}.
 \label{eq:symmetry}
\end{equation}
两个输入顺序共享模型和记录，平均值的方差取决于两次赋值的协方差。对称化后的概率质量应按这一读取规则重新评价。

互斥划分要求同组关系具有传递性。若 $A$ 与 $B$ 同组、$B$ 与 $C$ 同组，则 $A$ 与 $C$ 同组。三对判断出现两条正边和一条负边时，严格约束发生冲突。稀疏图还应检查正边连通分量内部的负边，以识别由较长路径产生的冲突。

概率也有整体相容要求。若 $p_{ij}$、$p_{jk}$ 和 $p_{ik}$ 拟解释为同一随机划分的边际概率，则共同条件信息下应满足
\begin{equation}
 p_{ik}\geq p_{ij}+p_{jk}-1.
 \label{eq:triangle}
\end{equation}
式\eqref{eq:triangle}来自同组事件的传递性及事件交集的概率下界。分别查询所得的关系概率可能与整体划分要求冲突，划分优化据此协调相互竞争的局部判断。专门模型的层级分类研究已观察到直接判断与分解判断之间的协调性差异\citep{joy2026}；聚类中的成对关系协调性应针对相应目标单独检验。

\subsection{从关系判断到整体划分}

令 $\mathcal P$ 表示记录集合的互斥划分，$u_{ij}(\mathcal P)=\ind\{i,j\text{ 同组}\}$。关系判断可以进入相关聚类中的加权不一致损失\citep{bansal2004}：
\begin{equation}
 \widehat{\mathcal P}=\argmin_{\mathcal P\in\mathfrak P}
 R_{\widehat p}(\mathcal P),\qquad
 R_p(\mathcal P)=\sum_{(i,j)\in E}a_{ij}
 \{p_{ij}(1-u_{ij})+(1-p_{ij})u_{ij}\}.
 \label{eq:partition}
\end{equation}
权重 $a_{ij}\geq0$ 且 $\sum_Ea_{ij}=1$，$\mathfrak P$ 为预先规定的划分空间。每一项是边际概率下的成对错误期望，期望的加总允许记录对之间存在相关性。算法对整体划分求解，所有同组指标共同满足传递性。若结合几何损失或簇数惩罚，其权重应在开发材料上确定。

概率 $p_{ij}$ 为同组判断提供连续权重。采用 $h_{ij}=\ind\{p_{ij}>1/2\}$ 时，式\eqref{eq:partition}相应最小化二元判断与划分之间的不一致数。进一步要求所有 $u_{ij}=h_{ij}$，即采用严格硬约束；其可行性取决于判断之间的冲突。二元不一致最小化和严格满足约束分别评价，便于观察误差处理的作用。

表\ref{tab:three}采用三条记录进行构造计算。设 $p_{AB}=0.9$、$p_{BC}=0.6$、$p_{AC}=0.2$，各边权重取一，报告损失总量。二元判断形成两条正边和一条负边，其不一致损失有三个并列最优划分；概率损失根据各边证据强度选择 $\{A,B\}\mid\{C\}$。

\begin{table}[htbp]
\centering\small
\caption{三条记录的关系冲突与整体划分：构造计算}\label{tab:three}
\begin{tabular}{@{}lrr@{}}
\toprule
划分 & 概率不一致损失 & 二元不一致数 \\
\midrule
$\{A,B,C\}$ & 1.3 & 1 \\
$\{A,B\}\mid\{C\}$ & 0.9 & 1 \\
$\{A,C\}\mid\{B\}$ & 2.3 & 3 \\
$\{B,C\}\mid\{A\}$ & 1.5 & 1 \\
$\{A\}\mid\{B\}\mid\{C\}$ & 1.7 & 2 \\
\bottomrule
\end{tabular}
\end{table}

\subsection{概率误差与查询覆盖}

关系获取的质量可以与整体损失建立明确联系。以下界来自式\eqref{eq:partition}的线性形式，证明见附录\ref{app:bound}。

\begin{proposition}\label{prop:bound}
固定查询集合 $E$、权重和划分空间。令 $r_{ij}$ 为目标关系的真实边际概率，$\mathcal P^*$ 为 $R_r$ 的最小化解。若计算所得 $\widehat{\mathcal P}$ 满足
$R_{\widehat p}(\widehat{\mathcal P})\leq\inf_{\mathcal P\in\mathfrak P}R_{\widehat p}(\mathcal P)+\varepsilon$，则
\begin{equation}
 R_r(\widehat{\mathcal P})-R_r(\mathcal P^*)
 \leq 2\sum_{(i,j)\in E}a_{ij}|\widehat p_{ij}-r_{ij}|+\varepsilon.
 \label{eq:bound}
\end{equation}
\end{proposition}

式\eqref{eq:bound}表明，查询集合 $E$ 上的超额成对风险受概率误差与数值求解误差约束。完整记录集合的划分质量还取决于查询覆盖。两个记录子集之间缺少关系信息时，其组间归属可能有多个同等支持的解。因此，稀疏查询下应分别计算已查询关系上的损失与完整划分的一致性。

查询设计应兼顾预算、困难记录和覆盖。近邻对可以提供细粒度区分，跨候选组的记录对可以检验桥接，随机对可以发现候选生成阶段遗漏的关系。基于表示选择候选图时，应在图外抽取评价对，检查同组关系的漏检。选择规则依赖前次判断时，应记录实际查询路径，并在来源比较中使用共同查询集合。

\subsection{评价聚类收益}

聚类评价需要参考划分与关系标准一致。调整兰德指数（ARI）评价划分的一致性，成对错误评价同组关系的恢复\citep{hubert1985}。语义编码的研究还应报告成对判断误差、概率评分、冲突率、簇数和不同运行的变化。内部几何指标可以描述所用表示中的紧密程度，目标机制划分则由按机制制定的参考材料评价。

研究可分别比较三种信息用途：直接使用同组关系、以三元组偏好训练表示，以及根据关系选择聚类粒度。三元组判断只比较锚点与两条记录的相对接近程度，其进入表示学习的方式不同于等价关系约束。固定下游算法和查询集合后比较关系来源，固定关系来源后比较二元与概率读取，能够分别评价信息来源及处理方式。第\ref{sec:numeric}节实施后一种比较。

\section{序数测量与记录链接：扩展应用}\label{sec:extensions}

\subsection{等级分布与测量变量}

序数测量依据有序规程记录程度。传统流程由评分者识别等级，语义决策模型可以依据等级描述提供分布 $\widehat p_{ik}$。Jev 的 Score 输出使用等级索引计算数值得分\citep{typesafe-score}。设研究者另选编码 $g(k)$，则相应汇总为
\begin{equation}
 \widehat s_i(g)=\sum_{k=0}^{K-1}g(k)\widehat p_{ik},\qquad
 \widehat t_i(c)=\sum_{k\geq c}\widehat p_{ik}.
 \label{eq:ordinal}
\end{equation}
前者依赖编码间距，后者表达达到给定等级阈值的赋值概率。序数回归利用等级次序建模\citep{mccullagh1980}；对数值均值的解释则需要对应的间距依据。

以三等级分布 $A=(0.6,0,0.4)$ 和 $B=(0,1,0)$ 为例。编码为 $(0,1,2)$ 时，得分分别为0.8和1；编码为 $(0,1,5)$ 时，得分分别为2和1。两种编码保持等级次序，却产生不同的得分排序。达到最高等级的概率始终为0.4和0。研究可以根据目标选择阈值比例、完整等级分布或具有实质间距依据的数值汇总。

评分评价应检查等级描述、评分一致性、极端等级的识别及不同群体的分布。\citet{gao2026}研究了决策模型对序数候选空间的使用，发现所研究模型最终选择的等级较参考标签更集中，提示评分研究应同时评价准确性与等级使用。对于跨组比较，规程应在各组表达相同程度，参考记录应覆盖各等级；多维构念可以分别评分，再按理论规定的关系进入统计模型。

\subsection{实体关系与链接后的统计单位}

记录链接识别来自不同记录的同一实体。传统概率链接依据字段比较证据评价匹配\citep{fellegi1969}；语义判断可以比较别名、简称及描述性信息。候选生成确定哪些记录对进入判断，语义比较确定其匹配证据，整体匹配规则确定最终实体关系。

设 $\widehat p_{ij}$ 为候选对的匹配概率，错误链接成本为 $C_{\mathrm{FP}}$，漏链接成本为 $C_{\mathrm{FN}}$。在概率对应匹配目标且采用这两项损失时，单对决策阈值为
\begin{equation}
 \widehat p_{ij}>\frac{C_{\mathrm{FP}}}{C_{\mathrm{FP}}+C_{\mathrm{FN}}}.
 \label{eq:linkage}
\end{equation}
两表一对一链接还须共同选择满足唯一匹配的记录对；实体去重则须形成具有传递性的同一实体关系。单对阈值与整体约束分别承担局部决策和结构组织。

链接评价宜同时报告候选覆盖、误链接与漏链接，以及实体数量、事件计数或回归结果的变化。错误合并可能减少实体数并混合属性，漏合并可能增加重复计数。语义判断因此参与统计单位的构造，其适用条件由字段证据、参考实体和分析用途共同确定。

\section{比较评价与研究议程}\label{sec:evaluation}

\subsection{构造聚类研究的设计}\label{sec:numeric}

本文以九条投诉原型开展受控数值研究，考察查询覆盖、判断可靠度与概率读取的作用。三种产品与三种主要故障机制完全交叉，每个组合对应一条记录，表\ref{tab:records}给出内容。参考划分依据主要机制形成三组。记录、参考关系和表示在重复计算中固定，变化来自查询集合及模拟判断。所有数值判断由构造机制产生，实验未调用通用语言模型或 Jev。

\begin{table}[htbp]
\centering\small
\caption{九条投诉原型：构造记录及预设主要机制}\label{tab:records}
\begin{tabularx}{\textwidth}{@{}L{1.9cm}YYY@{}}
\toprule
主要机制 & 智能门锁 & 烟雾报警器 & 室内温控器 \\
\midrule
供电模块失效 & 更换合格电池后仍无法通电，供电模块无输出 & 更换合格电池后仍无法通电，供电模块无输出 & 更换合格电池后仍无法通电，供电模块无输出 \\
断电后配置丢失 & 断电后已保存的密码丢失，需重新配置 & 断电后已保存的报警阈值丢失，需重新配置 & 断电后已保存的时段设置丢失，需重新配置 \\
传感测量失真 & 无人触碰时持续报告开门 & 室内无烟时持续报告烟雾 & 温度恒定时持续报告升温 \\
\bottomrule
\end{tabularx}
\end{table}

表示基线采用 $z_i=(2e_{\mathrm{product}(i)},e_{\mathrm{mechanism}(i)})$，其中 $e$ 为相应维度的单位向量。该表示明确设置产品维度较大的距离权重，用于展示目标机制与所用几何权重的差异。固定 $K=3$，枚举相应划分求得平方误差全局最小值6，其分组依据产品，与机制参考划分的 ARI 为 $-1/3$。这里使用构造表示；实际文本表示的维度权重和目标关系应由语料评价。

全部36个记录对构成候选图。每次随机排列这些对，前12、24和36对分别构成三个嵌套查询集合。每个条件重复1,000次，随机种子为20261005。关系方法在九条记录的全部21,147个划分上精确最小化式\eqref{eq:partition}，不预设簇数。并列解依次按与产品关系的不一致数、簇数及枚举顺序选择；这一规则在实验前固定，作为稀疏信息下的确定选择。参考机制标签用于构造和评价，未进入并列解选择。

设 $r_{ij}=\ind\{\mathrm{mechanism}(i)=\mathrm{mechanism}(j)\}$。均匀抽取记录对时，同组先验为 $\pi=9/36=0.25$。四种判断条件如下：

\begin{enumerate}
\item \textbf{完整正确判断。}查询返回 $p_{ij}=r_{ij}$，用于观察有限覆盖本身的作用。
\item \textbf{异质可靠度。}每个记录对以相同概率取 $h_{ij}=0.95$ 或0.80，模拟二元信号 $S_{ij}$ 以概率 $h_{ij}$ 等于 $r_{ij}$。返回概率按式\eqref{eq:posterior}计算。
\item \textbf{过度集中概率。}保留异质条件的全部信号和查询集合，计算返回概率时统一使用 $h=0.99$，考察高估可靠度的后果。
\item \textbf{标准偏移。}使用同一可靠度生成机制，但信号依据“产品相同”关系产生；统计评价仍采用机制参考划分，考察判断目标的偏移。
\end{enumerate}

异质条件的概率为
\begin{equation}
 p(S,h)=
 \begin{cases}
 \displaystyle\frac{h\pi}{h\pi+(1-h)(1-\pi)},&S=1,\\[0.7em]
 \displaystyle\frac{(1-h)\pi}{(1-h)\pi+h(1-\pi)},&S=0.
 \end{cases}
 \label{eq:posterior}
\end{equation}
先验与可靠度均由构造给定。式\eqref{eq:posterior}在均匀抽取记录对及所设噪声机制下提供边际概率；异质可靠度条件的校准对象为机制同组关系，标准偏移条件则依据产品同组关系生成判断。整体划分的相容性由优化程序处理。实际模型的概率质量应由开发与核验材料评价。

各条件均比较二元不一致最小化与概率加权，两种方法共享每次实验的查询图、划分空间及并列规则。另行记录全部二元约束可被严格满足的重复比例、成对判断误差及 Brier 分数。表示基线与关系方案比较不同信息输入，二元与概率方案则固定关系来源和数值求解过程。蒙特卡洛标准误根据各次实验的指标计算，配对差异的标准误根据同次实验的指标差计算。

\subsection{数值结果及其解释}

\begin{table}[htbp]
\centering\small
\caption{查询预算、判断质量与划分结果：每个条件重复1,000次}\label{tab:simulation}
\begin{tabular}{@{}lrrrrrr@{}}
\toprule
判断条件 & $B$ & 判断错误率 & Brier & 二元 ARI & 概率 ARI & 严格可行（\%） \\
\midrule
完整正确 & 12 & 0.0000 & 0.0000 & 0.3568 & 0.3568 & 100.0 \\
 & 24 & 0.0000 & 0.0000 & 0.8646 & 0.8646 & 100.0 \\
 & 36 & 0.0000 & 0.0000 & 1.0000 & 1.0000 & 100.0 \\
\addlinespace
异质可靠度 & 12 & 0.1297 & 0.0913 & 0.1711 & 0.1921 & 62.1 \\
 & 24 & 0.1266 & 0.0891 & 0.5067 & 0.5733 & 9.0 \\
 & 36 & 0.1265 & 0.0897 & 0.7336 & 0.7717 & 1.3 \\
\addlinespace
过度集中 & 12 & 0.1297 & 0.1241 & 0.1711 & 0.1719 & 62.1 \\
 & 24 & 0.1266 & 0.1211 & 0.5067 & 0.5172 & 9.0 \\
 & 36 & 0.1265 & 0.1210 & 0.7336 & 0.7495 & 1.3 \\
\addlinespace
标准偏移 & 12 & 0.4970 & 0.3503 & -0.2332 & -0.2308 & 62.6 \\
 & 24 & 0.5005 & 0.3519 & -0.2601 & -0.2481 & 7.9 \\
 & 36 & 0.4985 & 0.3515 & -0.2947 & -0.2620 & 1.3 \\
\bottomrule
\end{tabular}
\vspace{0.3em}
\parbox{\textwidth}{\footnotesize 所有指标以机制关系或机制划分为参考。ARI 为重复实验的平均调整兰德指数。二元 ARI 使用允许违反判断的最小不一致解；严格可行比例单独记录全部二元约束具有共同划分的次数。异质可靠度与过度集中条件共享二元信号，因此判断错误率、二元解及严格可行比例相同。全部未舍入指标和蒙特卡洛标准误随稿提供。}
\end{table}

完整正确条件下，12和24次查询的平均 ARI 分别为0.3568和0.8646，完全恢复参考划分的比例分别为0.6\%和39.4\%；完整查询图对应唯一的参考划分（表\ref{tab:simulation}）。有限查询即使全部正确，也可能支持多个整体结构，稀疏图上的结果因而取决于未查询关系及并列最优解的选择。

为考察稀疏查询中的解选择，本文在检查主要结果后补充敏感性分析，依次按最少簇数及原始枚举顺序选择并列最优解。完整正确条件下，12和24次查询出现多个最优划分的比例分别为99.5\%和60.6\%；替代规则的平均 ARI 分别为0.4905和0.9896。完整查询图仍恢复相同的正确划分。有限覆盖时，多个划分可在已查询关系上取得相同损失，却在未查询关系上作出不同分组。

带噪条件下也存在这一问题。异质可靠度条件在完整查询图上，二元与概率损失出现并列最优的比例分别为54.4\%和31.2\%。最少簇数优先规则下，二元 ARI 为0.8469，概率 ARI 为0.7771，方法排序与主规则相反。式\eqref{eq:bound}约束成对损失，ARI 评价两个划分的同组结构一致性；优化目标与评价指标的差异，以及并列解选择，均会影响比较结果。因此，信息使用方式的比较应同时报告目标损失、划分指标和解选择规则。

采用主规则时，异质可靠度条件下三个预算的概率 ARI 均高于二元 ARI。预算为24时，平均差异为0.0666，配对蒙特卡洛标准误为0.0058；预算为36时，差异为0.0381，标准误为0.0051。已知可靠度提供了判断强弱信息，概率损失据此分配权重。过度集中条件保留相同信号，但36次查询下的 Brier 分数由0.0897升至0.1210，概率 ARI 为0.7495，相对二元解的增益较小。在这一固定比较中，概率质量影响了加权后的划分表现。

严格约束的可行性随带噪判断累积而下降。异质条件中，严格可行比例从12次查询时的62.1\%降至36次查询时的1.3\%。在完整图上，具有两条正边和一条负边的三元组约占14.7\%。允许局部违反的整体优化仍给出划分，概率加权进一步利用证据强度。严格可行的少数重复与全体重复分开记录，以保持比较对象一致。

标准偏移条件下，三个预算的概率 ARI 均为负值。判断依据产品关系生成，增加查询进一步巩固了按产品分组的结构。此处的规程偏离机制目标，判断误差因而应相对于机制关系计算；置信度、重复一致性和查询数描述的是判断过程的其他属性。

上述结果基于九条固定原型、均匀查询、指定噪声机制与精确划分求解。异质可靠度条件中的先验和可靠度由构造给定，真实模型的概率质量及费用须在相应任务中测量。这些设定明确了查询覆盖、判断质量与使用方式之间的关系，可供真实语料研究逐项检验。

\subsection{局部判断、整体结构与下游结果}

第\ref{sec:title-case}节在观测数据上比较任务监督前后的分类误差与概率质量，本节在构造机制下考察关系覆盖、噪声、概率读取及解选择对整体划分的作用。前者固定部署条件比较任务适配，后者固定关系来源比较使用方式。分类概率可供复核或加权使用，聚类结果则同时取决于查询图及整体优化；两类任务各按相应统计目标评价。

附录\ref{sec:sms}报告 SMS 分数辅助总体比例估计的结果。该案例中，辅助分数与标签的总体均值接近，记录层面的相关系数为$-0.0545$，简单随机抽样（SRS）标签均值在所比较核验方案中的均方误差最低。评价辅助分数的抽样效率时，应考察其与分析变量的关系及所用统计设计。标签编码、关系结构和总体汇总因此分别设置评价目标。


\subsection{判断来源、读取方式与总成本}\label{sec:agenda}

标题案例的 NF4 配对设计固定基础模型、部署条件和输出读取，比较任务监督前后的表现。两个零样本部署同时改变量化及推理程序，各项配置的作用可通过后续单因素比较评价。分类研究还可在相同目标标签预算下比较通用模型、决策模型和传统监督方案，分别记录参数选择、训练及人工核验投入。标签预算曲线用于评价少标签条件下的效率，相同基础模型的适配比较用于评价训练目标的贡献。

真实语料的聚类来源比较可以采用表\ref{tab:comparison}中的方案。首先确定同一批记录、同一关系标准与独立参考划分，再固定候选图和查询集合。常规表示聚类用于评价所用表示提供的信息，通用语言模型与专门决策模型用于比较关系获取。共同返回的二元判断承担主要来源比较；具备概率字段的系统再开展二元与概率读取的内部比较。

\begin{table}[htbp]
\centering\small
\caption{语义标准指导的聚类：真实语料比较设计}\label{tab:comparison}
\begin{tabularx}{\textwidth}{@{}L{2.7cm}YY@{}}
\toprule
待评价因素 & 固定条件 & 比较内容与指标 \\
\midrule
额外关系信息 & 语料、目标标准、参考材料 & 常规表示聚类与关系辅助方案；ARI、成对误差及费用 \\
关系判断来源 & 输入记录对、标准、查询图、划分算法 & 通用语言模型与决策模型；关系误差、划分质量及总投入 \\
概率读取方式 & 同一来源、同一输出、同一划分空间 & 二元损失与概率损失；Brier、配对 ARI 差异、冲突 \\
查询选择 & 同一来源、标准、预算和统计方法 & 随机查询、近邻与跨组覆盖；完整关系恢复及图外漏检 \\
标准稳定性 & 同一语料和模型版本 & 等义改写、输入顺序、重复运行；目标判断及划分变化 \\
\bottomrule
\end{tabularx}
\end{table}

参考划分宜由领域人员依据目标规程形成，记录分歧及裁决依据。开发材料用于选择标准、概率校准和算法参数，评价材料用于最终比较。存在共享记录的成对判断时，误差具有相关性，重抽样应以记录或语料为单位保留这种结构。评价对同时覆盖候选图内部和外部，以检查候选生成的遗漏。

已有实测研究分别评价了计算社会科学标注中的费用与任务差异\citep{ibrahim2026}、直接与分解判断之间的概率协调性\citep{joy2026}，以及序数等级空间的使用\citep{gao2026}。这些结果涉及局部标注、跨问题结构和评分操作。语义标准指导的聚类研究应进一步联系关系误差、概率权重与整体划分，并报告查询、标准制定及人工核验的总投入。

后续研究可检验三个问题。目标任务标签较少且标准清楚时，语义判断能否降低达到既定误差水平的总成本？候选图和关系来源相同时，概率质量的改善能否提高目标划分的一致性？参考关系、查询预算和算法固定时，专门模型能否较通用模型提供更有价值的判断？各项结果可按语料、预算、类别或关系类型报告。

\section{结论}\label{sec:conclusion}

语义决策模型为统计分析提供依据自然语言标准获取标签、关系与等级的实现方式。其方法作用发生在监督信息和分析变量的构造环节：分类利用类别判断组织监督，聚类利用关系判断建立整体划分，序数测量利用等级分布表达程度，记录链接利用实体关系确定统计单位。既有零样本分类、弱监督和约束聚类为这些路径提供了方法基础。

本文依据研究标准、信息获取、读取方式和结构求解的联系，明确了各环节的评价对象。标题案例中，部署条件相同的任务微调提高了宏平均 F1，减少漏检、增加误报，概率评分也有所改善；其应用收益须结合参考标签、错误代价与标签投入评价。构造聚类研究中，查询覆盖影响完整划分的恢复，判断可靠度与并列解选择共同影响概率加权的表现，判断噪声影响全部约束共同满足的可行性，标准偏移则影响划分与研究目标的一致性。SMS 补充案例进一步表明，语义分数的下游价值取决于其与分析变量的关系及所用统计设计。

现有文献已展示语言模型判断参与标注与聚类的可行性，专门决策接口提供了相应的结构化操作。准确率、概率质量与总成本优势应在规定任务及比较条件下评价。研究者可先确定统计分析所需的信息类型，再选择相应接口，并依据关系结构、参考材料和预算确定使用方式。

\appendix
\section{成对风险误差界的证明}\label{app:bound}

对任意划分 $\mathcal P$，$u_{ij}\in\{0,1\}$，因而
\begin{align}
 |R_{\widehat p}(\mathcal P)-R_r(\mathcal P)|
 &=\left|\sum_E a_{ij}(\widehat p_{ij}-r_{ij})(1-2u_{ij})\right|\\
 &\leq\sum_E a_{ij}|\widehat p_{ij}-r_{ij}|=:\delta.
\end{align}
由 $\widehat{\mathcal P}$ 的近似最优性可得
\begin{align}
 R_r(\widehat{\mathcal P})
 &\leq R_{\widehat p}(\widehat{\mathcal P})+\delta\\
 &\leq R_{\widehat p}(\mathcal P^*)+\varepsilon+\delta\\
 &\leq R_r(\mathcal P^*)+\varepsilon+2\delta.
\end{align}
证明中的 $E$、权重与划分空间均固定。所得上界以查询图上的边际概率误差和数值求解误差约束超额成对风险；图外关系的恢复还取决于查询覆盖及其他结构信息。

\section{下游补充案例：SMS 分数与总体比例}\label{sec:sms}

总体估计可以作为语义变量的下游用途。本文复用既有 SMS Spam Collection 的抽样设计结果\citep{sms-data}。下载档案包含5,574行记录和747个垃圾短信标签，403次重复文本出现各保留为一个记录行；固定总体目标为归档标签比例0.1340。预测器采用固定版本的自然语言推断模型 \nolinkurl{cross-encoder/nli-MiniLM2-L6-H768}，对“正常通信”和“未经请求的推广或欺诈性垃圾短信”两项描述的蕴含 logit 实施 softmax，所得垃圾短信分数 $f_i$ 用作辅助变量\citep{sms-model}。该案例评价既有 NLI 读取与抽样设计。

若全体分数已知，从固定语料总体 $U$ 中简单随机不放回核验 $n$ 行，固定系数差值估计为
\begin{equation}
 \widehat\mu_\lambda=\lambda\bar f_U+
 \frac1n\sum_{i\in V}(Y_i^{\mathrm{ref}}-\lambda f_i),\qquad
 \Var_V(\widehat\mu_\lambda)=
 \left(1-\frac nN\right)\frac{S^2_{Y-\lambda f}}n.
 \label{eq:sms}
\end{equation}
其中 $N$ 为总体行数，$S^2$ 为分母 $N-1$ 的有限总体方差。辅助变量通过残差方差影响效率，这一原理与模型辅助估计及预测辅助推断衔接\citep{sarndal1992,ppi2023}。

两个核验预算各重复10,000次，结果摘要见表\ref{tab:sms}。分数与标签的相关系数为 $-0.0545$；全体标签计算的 SRS 最优固定系数为 $-0.1463$，理论方差降幅约0.30\%。这组分数通过固定系数提高抽样效率的空间较小。

\begin{table}[htbp]
\centering\small
\caption{复用 SMS 抽样设计结果：MSE 乘以 $10^4$}\label{tab:sms}
\begin{tabular}{@{}lrr@{}}
\toprule
估计方案 & 核验预算200 & 核验预算400 \\
\midrule
SRS 标签均值 & 5.568 & 2.637 \\
SRS 差值估计，$\lambda=1$ & 6.604 & 3.113 \\
先导分层校正，$\lambda=1$ & 10.117 & 3.603 \\
先导调权差值估计 & 9.642 & 3.413 \\
\bottomrule
\end{tabular}
\vspace{0.3em}
\parbox{\textwidth}{\footnotesize 两个先导方案均使用80个预算内先导标签，并沿用同一分层、分配和主阶段样本；调权系数由先导标签估计，在主阶段固定。完整设计、偏差、区间覆盖率和复现记录见既有 SMS 材料。表中目标均为归档标签的固定总体比例。}
\end{table}

在所比较核验方案中，SRS 标签均值的 MSE 最低。先导调权降低了单位权重先导方案的 MSE，结果同时受弱关联、系数估计变异和先导投入影响。全体预测均值为0.13818，接近标签比例，而记录层面的相关性较弱。总体汇总误差与辅助变量的抽样效率因而应分别评价。本文以总体估计检验语义信息的下游用途，完整推导与设计见既有材料。

\section{复现条件与材料}\label{app:materials}

标题案例的原始实验材料来自 \citet{title-case} 的固定 Git 提交 \nolinkurl{1b0ce28d265454680b5b383bbb7ff0929dca688b}。本稿保存对齐的六项逐条预测、数据处理清单、来源校验值和未舍入重分析结果，目录为 \nolinkurl{research/title_case/}。运行 \nolinkurl{scripts/title_case_analysis.py} 可复算分类与概率指标，并执行10,000次配对标题组自助重抽样；种子为20261006。正确提取的非空标题组有568个，错误提取有395个；组内重复出现保持原记录权重，24条空标题规则记录在每次重抽样中固定加入。该分析以一次训练得到的固定预测为条件。

Tev1 基础权重版本为 \nolinkurl{0b7becf017daa0e5eb222f8ce7483c8c8259c52f}。微调固定 NF4 四位主干、双重量化与 BF16 计算，LoRA 秩8、缩放系数16、dropout 0，学习率$5\times10^{-5}$，有效批量16，一轮训练，最长序列512词元。训练损失覆盖答案字母及结束标记；训练与验证提示均无截断，正式评价逐条推理。未微调对照在相同模型实例中关闭适配器。Ollama 0.35.0 使用 Q4\_K\_M GGUF，其 digest 随材料保存；该记录中的上游 HF 版本未标明。原实验使用 Python 3.12.11，NF4 推理与训练依赖 PyTorch 2.11.0、Transformers 5.14.1、PEFT 0.21.2 与 bitsandbytes 0.50.2，完整环境见来源仓库。此次重分析使用 Python 与 NumPy，沿用保存的预测。

聚类程序 \nolinkurl{scripts/semantic_clustering_study.py} 使用 Python 与 NumPy。程序通过受限增长串枚举全部划分，保存每次 ARI、严格约束可行性、汇总指标及蒙特卡洛标准误。保存目录为 \nolinkurl{research/semantic_clustering/}。本文报告的重复实验固定九条原型，随机变化仅来自记录对选择和判断噪声。

SMS 部分复用 \nolinkurl{research/real_data/results.json} 与 \nolinkurl{summary.csv}。其模型版本固定为 \nolinkurl{b95119ce93d3e065de6214e38cd4a97b0f2f2c6d}，预测文件校验值及读取描述见该目录的 \nolinkurl{provenance.json}。已有复现程序为 \nolinkurl{scripts/sms_design_evaluation.py}。语料总体行数依据原始下载档案解析结果报告。

本文的框架草拟、程序编写及文字修订使用生成式人工智能辅助。来源核查记录、构造设定及未舍入结果随稿提供。

\clearpage
\begin{thebibliography}{99}
\small
\interlinepenalty=10000

\bibitem[Almeida and Hidalgo(2011)]{sms-data}
Almeida, T. and Hidalgo, J. (2011).
\href{https://doi.org/10.24432/C5CC84}{SMS Spam Collection}.
UCI Machine Learning Repository. Dataset, accessed 5 October 2026.

\bibitem[Angelopoulos et~al.(2023)]{ppi2023}
Angelopoulos, A.~N., Bates, S., Fannjiang, C., Jordan, M.~I., and Zrnic, T. (2023).
\href{https://doi.org/10.1126/science.adi6000}{Prediction-powered inference}.
\emph{Science}, 382(6671), 669--674.

\bibitem[Bansal et~al.(2004)]{bansal2004}
Bansal, N., Blum, A., and Chawla, S. (2004).
\href{https://doi.org/10.1023/B:MACH.0000033116.57574.95}{Correlation clustering}.
\emph{Machine Learning}, 56, 89--113.

\bibitem[Dettmers et~al.(2023)]{qlora2023}
Dettmers, T., Pagnoni, A., Holtzman, A., and Zettlemoyer, L. (2023).
\href{https://proceedings.neurips.cc/paper_files/paper/2023/hash/1feb87871436031bdc0f2beaa62a049b-Abstract-Conference.html}{QLoRA: Efficient finetuning of quantized LLMs}.
\emph{Advances in Neural Information Processing Systems}, 36.

\bibitem[Fellegi and Sunter(1969)]{fellegi1969}
Fellegi, I.~P. and Sunter, A.~B. (1969).
\href{https://doi.org/10.1080/01621459.1969.10501049}{A theory for record linkage}.
\emph{Journal of the American Statistical Association}, 64(328), 1183--1210.

\bibitem[Gao et~al.(2026)]{gao2026}
Gao, T., Li, J., Li, Z., Chang, Y., and Wu, Y. (2026).
\href{https://arxiv.org/abs/2609.38827v1}{More choices, fewer decisions: Ordinal-scale bias in JEV-like direct-decision models}.
arXiv:2609.38827v1. Preprint.

\bibitem[Grimmer and Stewart(2013)]{grimmer2013}
Grimmer, J. and Stewart, B.~M. (2013).
\href{https://doi.org/10.1093/pan/mps028}{Text as data: The promise and pitfalls of automatic content analysis methods for political texts}.
\emph{Political Analysis}, 21(3), 267--297.

\bibitem[Guo et~al.(2017)]{guo2017}
Guo, C., Pleiss, G., Sun, Y., and Weinberger, K.~Q. (2017).
\href{https://proceedings.mlr.press/v70/guo17a.html}{On calibration of modern neural networks}.
\emph{Proceedings of Machine Learning Research}, 70, 1321--1330.

\bibitem[Hubert and Arabie(1985)]{hubert1985}
Hubert, L. and Arabie, P. (1985).
\href{https://doi.org/10.1007/BF01908075}{Comparing partitions}.
\emph{Journal of Classification}, 2, 193--218.

\bibitem[Ibrahim and Zaki(2026)]{ibrahim2026}
Ibrahim, H. and Zaki, Y. (2026).
\href{https://arxiv.org/abs/2609.24574v2}{Evaluating decision models for text annotation in computational social science}.
arXiv:2609.24574v2. Preprint, revised 23 September 2026.

\bibitem[Jia et~al.(2026)]{jia2026}
Jia, C., Wu, W., Guo, L., Lu, Z., Chen, C., and Ong, K.-L. (2026).
\href{https://arxiv.org/abs/2601.11118v1}{Optimized algorithms for text clustering with LLM-generated constraints}.
arXiv:2601.11118v1. AAAI-26 paper version.

\bibitem[Joy(2026)]{joy2026}
Joy, S.~S. (2026).
\href{https://arxiv.org/abs/2609.33971v1}{Do System One decisions add up? A study of probabilistic coherence}.
arXiv:2609.33971v1. Preprint.

\bibitem[LeoAndJellyfish(2026)]{title-case}
LeoAndJellyfish (2026).
\href{https://github.com/LeoAndJellyfish/paper-title-error-classification/tree/1b0ce28d265454680b5b383bbb7ff0929dca688b}{论文标题提取错误识别：课程实验与复现材料}.
GitHub repository, fixed commit specified in Appendix~\ref{app:materials}, accessed 6 October 2026.

\bibitem[McCullagh(1980)]{mccullagh1980}
McCullagh, P. (1980).
\href{https://doi.org/10.1111/j.2517-6161.1980.tb01109.x}{Regression models for ordinal data}.
\emph{Journal of the Royal Statistical Society, Series B}, 42(2), 109--127.

\bibitem[Ratner et~al.(2017)]{ratner2017}
Ratner, A., Bach, S.~H., Ehrenberg, H., Fries, J., Wu, S., and R{\'e}, C. (2017).
\href{https://doi.org/10.14778/3157794.3157797}{Snorkel: Rapid training data creation with weak supervision}.
\emph{Proceedings of the VLDB Endowment}, 11(3), 269--282.

\bibitem[S{\"a}rndal et~al.(1992)]{sarndal1992}
S{\"a}rndal, C.-E., Swensson, B., and Wretman, J. (1992).
\href{https://link.springer.com/book/9780387406206}{\emph{Model Assisted Survey Sampling}}.
Springer, New York.

\bibitem[Sentence Transformers(n.d.)]{sms-model}
Sentence Transformers (n.d.).
\href{https://huggingface.co/cross-encoder/nli-MiniLM2-L6-H768}{Cross-encoder for natural language inference: nli-MiniLM2-L6-H768}.
Model card and fixed model files; revision specified in Appendix~\ref{app:materials}, accessed 5 October 2026.

\bibitem[Tan et~al.(2024)]{tan2024}
Tan, Z., Li, D., Wang, S., Beigi, A., Jiang, B., Bhattacharjee, A., Karami, M., Li, J., Cheng, L., and Liu, H. (2024).
\href{https://doi.org/10.18653/v1/2024.emnlp-main.54}{Large language models for data annotation and synthesis: A survey}.
In \emph{Proceedings of EMNLP}, pp.~930--957.

\bibitem[Together AI(2026)]{tev1-card}
Together AI (2026).
\href{https://huggingface.co/togethercomputer/Tev1-4B-experimental/tree/0b7becf017daa0e5eb222f8ce7483c8c8259c52f}{Tev1-4B-experimental}.
Model card, fixed revision specified in Appendix~\ref{app:materials}, accessed 6 October 2026.

\bibitem[TypeSafe AI(2026a)]{typesafe-choice}
TypeSafe AI (2026a).
\href{https://docs.typesafe.ai/primitives/choice}{Choice}.
Technical documentation, accessed 5 October 2026.

\bibitem[TypeSafe AI(2026b)]{typesafe-intro}
TypeSafe AI (2026b).
\href{https://docs.typesafe.ai/introduction}{Introduction}.
Technical documentation, accessed 5 October 2026.

\bibitem[TypeSafe AI(2026c)]{typesafe-noul}
TypeSafe AI (2026c).
\href{https://docs.typesafe.ai/primitives/noul}{Noul}.
Technical documentation, accessed 5 October 2026; response definition rechecked 5 October 2026.

\bibitem[TypeSafe AI(2026d)]{typesafe-score}
TypeSafe AI (2026d).
\href{https://docs.typesafe.ai/primitives/score}{Score}.
Technical documentation, accessed 5 October 2026.

\bibitem[Viswanathan et~al.(2024)]{viswanathan2024}
Viswanathan, V., Gashteovski, K., Lawrence, C., Wu, T., and Neubig, G. (2024).
\href{https://doi.org/10.1162/tacl_a_00648}{Large language models enable few-shot clustering}.
\emph{Transactions of the Association for Computational Linguistics}, 12, 321--333.

\bibitem[Wagstaff et~al.(2001)]{wagstaff2001}
Wagstaff, K., Cardie, C., Rogers, S., and Schroedl, S. (2001).
\href{https://www.cs.cmu.edu/~dgovinda/pdf/icml-2001.pdf}{Constrained k-means clustering with background knowledge}.
In \emph{Proceedings of ICML}, pp.~577--584.

\bibitem[Wang et~al.(2025)]{cequel2025}
Wang, H., Zhang, T., Yang, R., and Xu, J. (2025).
\href{https://arxiv.org/abs/2504.15640v2}{Cequel: Cost-effective querying of large language models for text clustering}.
arXiv:2504.15640v2. Preprint, revised 21 August 2025.

\bibitem[Yin et~al.(2019)]{yin2019}
Yin, W., Hay, J., and Roth, D. (2019).
\href{https://doi.org/10.18653/v1/D19-1404}{Benchmarking zero-shot text classification: Datasets, evaluation and entailment approach}.
In \emph{Proceedings of EMNLP-IJCNLP}, pp.~3914--3923.

\bibitem[Zhang et~al.(2023)]{clusterllm2023}
Zhang, Y., Wang, Z., and Shang, J. (2023).
\href{https://doi.org/10.18653/v1/2023.emnlp-main.858}{ClusterLLM: Large language models as a guide for text clustering}.
In \emph{Proceedings of EMNLP}, pp.~13903--13920.

\end{thebibliography}
\end{document}
